% Part 3: more useful concepts. Divider plus 3 frames, session 2:08-2:36.
\section{Part 3. More useful concepts}\subsection{}

% Part 3 divider.
\begin{transitionframe}
Part 3. More useful concepts
\end{transitionframe}

% The eight working tips, as one list.
\begin{frame}{Working tips}
\begin{enumerate}
    \item Plan before acting.
    \item Be specific.
    \item Always have a backup.
    \item Manage your context window.
    \item Always have a way to verify outputs.
    \item Ask the LLM to take notes.
    \item Be replicable.
    \item Stay organized with your code, data, and notes.
\end{enumerate}
\end{frame}

% Permissions, honestly.
\begin{frame}{AI safety: honestly, you will grant most permissions}
\begin{wideitemize}
    \item Every step today asked for file writes, shell access and network access. It is unrealistic to approve every single permission every time.
    \item The best you can do is to allow-list broadly and stop reading.
    \item Which is tip 3 again, for a different reason: \ybb{always have a backup.}
    \item Mine is GitHub for code and Dropbox for data, and my own agents run on \ybb{a separate Mac Mini} -- a machine with limited access to my data that I can wipe.
    \item The real question is not ``what may it do'' but \yob{``what can it reach''}.
\end{wideitemize}
\end{frame}

% Dev containers: what isolation buys and what it does not.
\begin{frame}{Dev containers are the real alternative}
\begin{wideitemize}
    \item Cost: Docker, a config file, and a slower first run. Windows needs WSL2.
    \item Isolation is about \ybb{what is reachable}. 
    \item See ``\#7 Permissions, Sandboxes, and Autonomous Agents'' in \emph{A Series on Claude Code}, by Paul Goldsmith-Pinkham.
\end{wideitemize}
\end{frame}
